% !PRINTED 793
\weberpage{793}
\begin{flushleft}
Symmetrische Gruppe, ihre Normaltheiler I, 602. \\
Syzygetische und azygetische Systeme von Doppeltangenten II, 362. \\
— — — Complexe von Doppeltangenten II, 366. \\[6pt]
\textbf{T}angenten einer Curve II, 326. \\
\gesperrt{Taylor}'sche Entwickelung I, 59. \\
Ternäre Formen II, 322. \\
— lineare Gruppen vom Grade 168 II, 433. \\
— — Congruenzgruppe v. Grade 168 II, 475. \\
Tetra\"edergruppe II, 212. \\
— im Congruenzk\"orper II, 276. \\
Theilbarkeit ganzer Functionen I, 32. \\
— ganzer Functionale II, 509. \\
— ganzer Zahlen II, 510. \\
Theiler von ganzen Zahlen I, 1. \\
— von ganzen Zahlen, gr\"osster gemeinschaftlicher I, 359. \\
— ganzer Functionen I, 27; II, 497. \\
— ganzer Functionen, gr\"osster gemeinschaftlicher I, 34. \\
— einer Gruppe I, 476, 501; II, 7. \\
— der Ikosa\"edergruppe II, 227. \\
— eines K\"orpers I, 450. \\
— der linearen Congruenzgruppe II, 271. \\
Theilerfremde Zahlen I, 2, 359. \\
— Functionen I, 35. \\
— Functionale II, 514. \\
Theilnenner eines Kettenbruches I, 361. \\
Theilung des Winkels I, 435. \\
Totalresolventen I, 512. \\
Tr\"agheitsgesetz quadratischer Formen I, 183, 255. \\
Transcendente Functionen zur L\"osung der Ikosa\"edergleichung II, 429. \\
— Zahlen II, 745. \\
Transcendenz der Zahl $e$ II, 751. \\
— der Zahl $\pi$ II, 756. \\
Transformation einer Gruppe I, 506. \\
— der quadratischen Form in eine Summe von Quadraten I, 181. \\
— von Formen $n^{\text{ten}}$ Grades I, 185. \\
— der cubischen Gleichung I, 219. \\
Transformation der Gleichung f\"unften Grades I, 230. \\
Transformirte Substitutionen II, 156. \\
Transitive und intransitive Gruppen I, 482, 506. \\
— Permutationsgruppe vom Index 6 von sechs Ziffern I, 628. \\
Transponirte Substitutionen II, 156. \\
Transpositionen I, 65, 492. \\
Trigonometrische L\"osung reiner Gleichungen I, 111. \\
— — cubischer Gleichungen I, 349. \\
Trinomische Gleichungen I, 352. \\
Tripelgleichungen II, 342. \\
Tripelsysteme II, 310. \\
\gesperrt{Tschirnhausen}-Transformation I, 170, 210. \\
— — der cubischen Gleichung I, 219. \\[6pt]
\textbf{U}mkehrbare Perioden I, 391. \\
Unbekannte I, 20. \\
Unbestimmte Gleichungen I, 365. \\
Unendliche Wurzeln einer Gleichung I, 136. \\
Unendlichkeit ganzer rationaler Functionen I, 104. \\
Unicursalcurven II, 416. \\
Unterdeterminanten I, 72. \\
— h\"ohere I, 77. \\
— complement\"are I, 79. \\
Unz\"ahlbare Mengen II, 748. \\
Urspr\"ungliche Functionen I, 25. \\[6pt]
\textbf{V}ariable in der K\"orpertheorie II, 499. \\
Verh\"altnisse I, 10. \\
Verwandte Zahlen und Functionale II, 631. \\
Volles Restsystem I, 359; II, 541. \\
Vollst\"andige Systeme von Doppeltangenten II, 367. \\
Volumen II, 561, 690. \\
Vorzeichenwechsel ganzer Functionen I, 107. \\[6pt]
\textbf{W}endepunkte II, 326. \\
Wendetangenten II, 326. \\
— einer Curve dritter Ordnung II, 331. \\
\gesperrt{Wilson}'scher Lehrsatz I, 428. \\
Winkeltheilung I, 551, 609.
\end{flushleft}
